\subsection{A common Frobenius polynomial for the actual representations}
\label{sec:support-spectrum}

Fix an imaginary quadratic field $K$, and put $d_K=|\operatorname{disc}K|$.
Constants in this subsection may depend on $K$.  They are independent of
the conductor, the twist, both auxiliary primes, every coefficient prime,
and every coefficient field.  This order of quantifiers is essential.

Here is the precise local condition used below.  If $F$ is a number field,
$M\ge1$, and $r$ is an $n$-dimensional $\ell$-adic representation, say
that $r$ has the original level condition with exponents
$m_1>\cdots>m_n$ and modulus $M$ if it is continuous and has the
following properties.  At each finite $v\nmid\ell$, its restriction
to inertia is diagonal with characters that factor through
$\operatorname{Gal}(F(\mu_M)/F)$.  At each $v\mid\ell$, it has a
full invariant flag whose successive characters $\kappa_i$ satisfy
\[
 \kappa_i|_{I_v}=\varepsilon^{m_i}\chi_i|_{I_v},\qquad
 \chi_i\text{ factors through }\operatorname{Gal}(F(\mu_M)/F).
                                                               \tag{SP1}
\]
In addition, either $r|_{G_{F_v}}$ is potentially crystalline, or
$\kappa_i\ne\kappa_j\varepsilon$ for every $i,j$.  The latter
condition compares the whole characters of $G_{F_v}$, not only
their restrictions to inertia.  This is exactly the original
local condition, including its alternative at the coefficient
prime \cite{support-spectrum-source}.

We use two classes of representations, with the source's normalizations.
In rank three, $r$ is absolutely irreducible, has (SP1) with
$(m_1,m_2,m_3)=(e,0,-e)$ and $M=N$, and
$r\varepsilon^{-1}$ is attached to a regular algebraic,
conjugate self-dual cuspidal representation $\pi$ of
$\mathrm{GL}_3(\mathbf A_K)$.  Thus, at an unramified place
$w$ of residue cardinality $Q$, if $\beta_{\pi,w,i}$ are the
unitarily normalized Satake parameters, then
\[
 \operatorname{Eig}\bigl((r\varepsilon^{-1})
                         (\operatorname{Frob}_w^{-1})\bigr)
       =\{Q\beta_{\pi,w,1},Q\beta_{\pi,w,2},Q\beta_{\pi,w,3}\}.
                                                               \tag{SP2}
\]
Throughout, $\operatorname{Frob}_w$ is arithmetic Frobenius,
so $\varepsilon(\operatorname{Frob}_w^{-1})=Q^{-1}$.
In rank two, $r\simeq W|_{G_K}\oplus1$, where $W$ is an
irreducible representation of $G_{\mathbf Q}$ with (SP1) for
$(e,-e)$ and $M=Nd_K$, and $W\varepsilon^{-e}$ is attached to
a holomorphic normalized newform of weight $2e+1$.  The actual
rank-two class also has the original weight-zero determinant
condition.  The counting argument below holds for the larger class
just described, so requires no change to that condition or to $W$.
Theorems~\ref{thm:support-type-one-rigidity}
and~\ref{thm:support-type-two-rigidity} place the actual
representations from the integral Selmer model in these classes,
with the indicated Tate twists and the original level bounds.

\begin{theorem}[The original uniform Frobenius polynomial]
\label{thm:support-spectrum}
There are positive integers $C_K,A_K$ with the following property.
Let $e,N\ge1$, let $p\ge7$ be a rational prime not dividing
$Nd_K$, and let $w\mid p$, with $Q=\#(\mathcal O_K/w)$.
There is a finite multiset of algebraic numbers
$\{\alpha_s:s\in\Sigma\}$, independent of the coefficient prime,
such that
\[
 \begin{split}
 R:=|\Sigma|&\le C_K[N(e+1)]^{A_K},\\
 |\sigma(\alpha_s)|&\le3
       \quad\text{for every complex embedding }\sigma,\\
 Q^e\alpha_s&\in\overline{\mathbf Z}.
 \end{split}                                                   \tag{SP3}
\]
The indexed multiset is stable under
$\operatorname{Gal}(\overline{\mathbf Q}/\mathbf Q)$.
For every $\ell\ne2,p$, every isomorphism
$\iota:\overline{\mathbf Q}_\ell\simeq\mathbf C$, and every
finite family $X$ of pairwise nonisomorphic representations of the
two original classes, there is an injection $j:X\hookrightarrow\Sigma$
such that
\[
 \iota\!\left(\operatorname{tr}r(\operatorname{Frob}_w^{-1})\right)
                 =\alpha_{j(r)}.                              \tag{SP4}
\]
Repeated trace values retain their multiplicities in $\Sigma$.

Consequently, for every finite Galois-stable indexed multiset
of roots of unity $\{z_o:o\in\mathcal O\}$, there is an integer
$B\ne0$, independent of $\ell$ and $\iota$, satisfying
\[
 \log|B|\le C_K|\mathcal O|[N(e+1)]^{A_K}\log p.                \tag{SP5}
\]
After increasing the same constants if necessary, for every family
$X_o$ as above one has
\[
 \begin{split}
 \operatorname{tr}r(\operatorname{Frob}_w^{-1})
   &\ne\iota^{-1}(Q^{-e}z_o+1+Q^ez_o^{-1}),\\
 \sum_o\sum_{r\in X_o}v_\ell\!\left(
   \operatorname{tr}r(\operatorname{Frob}_w^{-1})
      -\iota^{-1}(Q^{-e}z_o+1+Q^ez_o^{-1})\right)
   &\le v_\ell(B),\qquad v_\ell(\ell)=1.
 \end{split}                                                   \tag{SP6}
\]
These are the full original Frobenius congruence-number requirements.
\end{theorem}

\begin{proof}
We first bound the level, including its factor at the coefficient
prime.  We then count a common spectrum, check its integral
normalization, and form one integer.

\emph{The coefficient-prime flag.}
Let $F_v/\mathbf Q_\ell$ be the local field.  A character
in (SP1) is de Rham, since after a finite extension it is an
unramified character times a power of $\varepsilon$.  If the
potentially crystalline alternative in (SP1) holds, its Weil--Deligne
monodromy is zero and its inertia characters are the $\chi_i$.
Suppose instead that all the stated character inequalities hold.
We prove the same assertion from the flag.

For a de Rham representation $V$ all of whose Hodge--Tate weights
are strictly negative, Bloch--Kato's dimension formula and local
Euler characteristic give
\[
 \begin{split}
 \dim H_f^1(F_v,V)&=[F_v:\mathbf Q_\ell]\dim V+\dim H^0(F_v,V),\\
 \dim H^1(F_v,V)&=[F_v:\mathbf Q_\ell]\dim V+\dim H^0(F_v,V)
                          +\dim H^0(F_v,V^*(1)).
 \end{split}                                                   \tag{SP7}
\]
The first formula applies to de Rham representations, not only
to representations already crystalline over $F_v$
\cite[Corollary~3.8.4]{bloch-kato-tamagawa}.
Our convention is that $\varepsilon$ has Hodge--Tate weight $-1$.
For the extension of $\kappa_2$ by $\kappa_1$, the relevant
representation is $\kappa_1\kappa_2^{-1}$.  Its cyclotomic
exponent is $e$ in rank three and $2e$ in rank two, so its
Hodge--Tate weight is strictly negative.  The nonresonance
assumption is precisely
$H^0(F_v,\kappa_1^{-1}\kappa_2\varepsilon)=0$.
Thus (SP7) gives $H^1=H_f^1$.

In rank three, let $U$ be this two-dimensional extension.
It is de Rham.  The final extension is measured by
$V=U\kappa_3^{-1}$.  Its successive cyclotomic exponents
are $2e,e$, hence both Hodge--Tate weights are negative;
the two possible constituents of $V^*(1)$ are nontrivial
by the same original nonresonance condition.  Therefore
$H^0(F_v,V^*(1))=0$ and again $H^1=H_f^1$.
After a finite extension killing the $\chi_i$, these finite
extension classes are crystalline extension classes.  Inducting
up the flag proves potential crystallinity.  Exactness of the
potentially semistable comparison functor shows that its finite
inertia representation has the successive characters $\chi_i$;
finite inertia is semisimple in characteristic zero.  Its
monodromy is zero.  The possibly large degree of this finite
extension enters none of the bounds below.

This argument includes $e=1$.  In rank three that is exactly
where a ratio can equal $\varepsilon$ and give a noncrystalline
semistable extension.  We used the explicit alternative in
(SP1) to exclude that case; we did not infer crystallinity from
the weights alone.

\emph{The conductor at every finite place.}
For a finite character factoring through
$\operatorname{Gal}(F(\mu_M)/F)$, local reciprocity shows
that its conductor exponent at $v$ is at most $v(M)$.
Indeed the character extends from the image of $G_F$ in
$(\mathbf Z/M\mathbf Z)^\times$ to a character of that
finite abelian group.  If $v\mid M$, the norm of an element
of $1+M\mathcal O_{F_v}$ is in $1+M\mathbf Z_\ell$:
write multiplication by the element in an integral basis and
take its determinant.  Norm compatibility in local reciprocity
therefore kills the character on $1+M\mathcal O_{F_v}$.
If $v\nmid M$, the cyclotomic extension in question is
unramified at $v$, giving conductor exponent zero.

At $v\nmid\ell$, (SP1) already makes inertia finite and
diagonal, so the monodromy is zero.  At $v\mid\ell$, this
was proved above.  In both cases additivity of the Artin
conductor gives
\[
       a_v(\operatorname{WD}(r))\le n\,v(M).                  \tag{SP8}
\]
Tate twists are unramified on Weil--Deligne inertia and do
not change this conductor.

The original attachment condition gives equality of characteristic
polynomials at almost all Frobenius elements.  Chebotarev and
semisimple character theory therefore identify the original
$r\varepsilon^{-1}$, or $W\varepsilon^{-e}$, with the
corresponding compatible automorphic Galois representation.
Full local--global compatibility consequently identifies these
Weil--Deligne representations, including their monodromy,
with the local automorphic parameters.  For the rank-three
representations use Caraiani's Theorem~1.1 both away from
and at the coefficient prime; for classical newforms use
the classical local compatibility theorem and Saito's
coefficient-prime comparison with its weight-monodromy supplement
\cite{support-spectrum-Caraiani2012,support-spectrum-Caraiani2014,
support-spectrum-Saito}.  These theorems have no bound on
the weight relative to $\ell$.  Equality only after ordinary
semisimplification would not justify (SP8) for the
automorphic conductor.  The local new-vector theorem,
with its corrected proof, now gives fixed vectors under
\[
 K_3(N^3)=\ker\!\left(\mathrm{GL}_3(\widehat{\mathcal O}_K)
       \longrightarrow\mathrm{GL}_3(\mathcal O_K/N^3)\right)
                                                               \tag{SP9}
\]
in rank three, and a classical newform level dividing $M^2$
in rank two \cite{support-spectrum-newvector}.  Passing to
these larger counting levels keeps every original representation
and its polarization; it changes neither its local condition
nor its trace.

\emph{Counting the common spectra.}
The local Hodge--Tate weights in rank three determine the
infinitesimal character of $\pi$ to be $(e,0,-e)$ at both
complex embeddings.  The corresponding algebraic coefficient
representation of $\operatorname{Res}_{K/\mathbf Q}
\mathrm{GL}_3$ has highest weight
\[
                      (e-1,0,1-e)                             \tag{SP10}
\]
at each embedding.  The Weyl dimension formula gives dimension
$e^3$ at either embedding and $e^6$ for the full coefficient
system.  This remains valid at $e=1$, when that coefficient
system is trivial.  The infinitesimal character is still regular.

Let $\Sigma_3$ be the set of regular algebraic conjugate self-dual
cuspidal $\pi$ with (SP10) and a $K_3(N^3)$-fixed vector.
Fix once a neat compact open subgroup $U_0$ of
$\mathrm{GL}_3(\widehat{\mathcal O}_K)$, and let
$U=U_0\cap K_3(N^3)$.  Borel--Serre compactification gives
a finite cell complex for the locally symmetric space at
level $U_0$, with finitely many components.  The space at
level $U$ is a finite covering of degree at most
\[
 [U_0:U]\le
 |\mathrm{GL}_3(\mathcal O_K/N^3)|\le N^{54}.                  \tag{SP11}
\]
The split real centre is divided out in defining these spaces.
The coefficient system (SP10) has trivial central weight.
Lift a fixed cell structure through the covering.  Cellular
cochains with this local system have total dimension at most
the number of lifted cells times $e^6$.

Every member of $\Sigma_3$ contributes a nonzero summand to
cuspidal cohomology with these coefficients, and cuspidal
cohomology injects into ordinary cohomology.  The standard
cohomological decomposition and strong multiplicity one
separate the summands.  Borel--Serre and Clozel's
Lemmas~3.14--3.15 therefore give
\[
                  |\Sigma_3|\le C_KN^{54}e^6.                 \tag{SP12}
\]
This is a count of representations, so no upper bound for
their Hecke fields is needed
\cite{support-spectrum-BorelSerre,support-spectrum-Clozel}.
All dependence on the fixed field, including the number of
components and cells, is confined to $C_K$.

Let $\Sigma_2$ be the set of normalized primitive holomorphic
newforms of weight $k=2e+1$ and level dividing $M^2=(Nd_K)^2$,
with all nebentypes included.  Their eigenlines are independent
in the space of cusp forms at level $\Gamma_1(M^2)$.
For completeness a coarse uniform dimension bound suffices.
Intersect with a fixed neat principal subgroup of level three.
The covering degree is at most a constant times $M^8$, by
bounding it by the number of two-by-two matrices modulo $M^2$.
The corresponding modular curve has a cell decomposition lifted
from a fixed finite one, and the local system
$\operatorname{Sym}^{k-2}$ has rank $k-1=2e$.
The holomorphic Eichler--Shimura summand injects into its first
cohomology.  Hence
\[
                  |\Sigma_2|\le C e(Nd_K)^8.                 \tag{SP13}
\]
This also holds for the smallest weight $k=3$.

Clozel's rationality theorem makes $\Sigma_3$ Galois-stable
\cite[Theorem~3.13]{support-spectrum-Clozel}.
Scalar conjugation preserves both the fixed finite level and
conjugate self-duality.  It permutes the two coefficient
embeddings of $K$, whose weights in (SP10) agree.
The analogous statement for $\Sigma_2$ is the classical
rationality of normalized newforms and their Hecke action.
Thus both spectra are finite and Galois-stable.  Notice that
the rank-three normalization (SP2) uses the rational number
$Q$, not a square root of $Q$.

\emph{Traces and integrality.}
For $\pi\in\Sigma_3$, put
$\alpha_\pi=\sum_i\beta_{\pi,w,i}$.  Caraiani's temperedness
theorem gives $|\beta_{\pi,w,i}|=1$, hence
$|\alpha_\pi|\le3$; the same holds for every conjugate
\cite[Theorem~1.2]{support-spectrum-Caraiani2012}.
Rationality makes these numbers algebraic.  To check the
denominator, consider the compatible Galois representation
for $\pi$ at each coefficient prime, in the normalization
obtained by multiplying (SP2) by $\varepsilon$.
At any coefficient prime $\ell\ne p$, continuity gives a
stable integral lattice.  Since the local automorphic
representation at $w$ is unramified, full local compatibility
makes this Galois representation unramified there.  Its
Frobenius eigenvalues, and their inverses, are integral at $\ell$.

At the coefficient prime $p$, the representation is crystalline
at $w$, with Hodge--Tate weights $-e,0,e$ at each embedding.
The existence, Hodge--Tate description, and crystalline assertion
are the basic automorphic Galois-representation theorem
\cite[Theorem~2.1.1]{support-spectrum-BLGGT}; unramified local
components are within its Iwahori-fixed case.
Weak admissibility puts each Newton slope between $-e$ and $e$.
Equivalently, for the eigenvalues of $\varphi^f$, where
$Q=p^f$, their valuations normalized by $v_p(p)=1$ lie
between $-ef$ and $ef$.  For the upper bound, one can also
apply the lower bound to the dual representation.
These eigenvalues are the parameters $\beta_{\pi,w,i}$ in
the normalization (SP2).  Thus $Q^e\alpha_\pi$ is integral
at every place above $p$, as well as at every other finite
place.  It is an algebraic integer.

For $g\in\Sigma_2$, let $\lambda_1,\lambda_2$ be the roots of
\[
                 T^2-a_p(g)T+\eta_g(p)p^{2e}.                 \tag{SP14}
\]
Here $p\nmid M$, so these are the usual good-prime roots.
The modular-form Galois representation uses inverse arithmetic
Frobenius in (SP14), exactly as in the preceding rank-two
comparison.  Set
\[
       \alpha_g=1+p^{-ef}(\lambda_1^f+\lambda_2^f),
               \qquad f=[\mathcal O_K/w:\mathbf F_p]\in\{1,2\}.
                                                               \tag{SP15}
\]
Deligne's bound gives $|\lambda_i|=p^e$, so
$|\alpha_g|\le3$.  The roots in (SP14) are algebraic
integers, and (SP15) gives $Q^e\alpha_g$ integral directly.
The same facts hold for each Galois conjugate
\cite{support-spectrum-Deligne}.

Take $\Sigma$ to be the disjoint union $\Sigma_3\sqcup\Sigma_2$
with these traces.  The bounds (SP12)--(SP13) prove (SP3).
Given an original representation of the first class, its
attached $\pi$ belongs to $\Sigma_3$ by (SP8)--(SP10), and
(SP2) gives (SP4).  If two such representations had the same
$\pi$, their characteristic polynomials at almost all
Frobenius elements would agree; Chebotarev and semisimple
character theory would make them isomorphic.  In the second
class choose its primitive newform.  The same newform gives
the same $W$, hence the same $W|_{G_K}\oplus1$.
Thus a pairwise nonisomorphic family injects into $\Sigma$.
Different newforms can restrict to isomorphic representations
over $K$, but this only leaves unused indices in $\Sigma_2$;
it never invalidates the injection.  Repeated numerical
traces are not removed.

\emph{One integer for all coefficient primes.}
Put $a=Q^e$ and form
\[
 \begin{gathered}
 P(T)=\prod_{s\in\Sigma}(T-a\alpha_s)\in\mathbf Z[T],\qquad
 Z(X)=\prod_{o\in\mathcal O}(X-z_o)\in\mathbf Z[X],\\
 B=\operatorname{Res}_X\!\left(Z(X),X^R P(X+a+a^2/X)\right).
 \end{gathered}
                                                               \tag{SP16}
\]
The verified spectra now satisfy every hypothesis of
Proposition~\ref{selmer-resultant-certificate}.  It gives
$B\ne0$, all the trace separations in (SP6), and
\[
             \log|B|\le |\mathcal O|R(2e\log Q+\log2).        \tag{SP17}
\]
Since $f\le2$, the extra factor $e$ and the factor $\log Q$
are absorbed in (SP5) by increasing the polynomial exponent
once.  For example the coarse bounds above allow $A_K=55$;
the precise exponent has no role below.

The injections (SP4), separately for each $o$, select
subsets of its indexed factors, including their multiplicities.
The valuation conclusion of that proposition gives (SP6).
Empty spectra or empty multisets give $B=1$ and satisfy
the same assertions.
\end{proof}

For clarity, the good-data hypotheses in the next corollary are
the following explicit ones.  Write
$\psi=\psi^{(\ell)}\psi_\ell$ for its decomposition by the
prime-to-$\ell$ and $\ell$-power parts of the modulus, and let
$F_{\psi^{(\ell)}}$ be the abelian field cut out by
$\psi^{(\ell)}$.  They are
\[
 \begin{gathered}
 \ell\ne2,\quad e\ge1,\quad 2\mid N,\quad
 \ell\nmid\operatorname{ord}\psi,\quad \psi_\ell^2=1,
 \quad\psi(-1)=(-1)^e,\\
 K\text{ imaginary quadratic},\quad
 \ell\nmid d_K h_K,\quad (d_K,N)=1,\quad
 K\not\subset\mathbf Q(\mu_\ell)F_{\psi^{(\ell)}},\quad
 \ell-1\nmid2e,
 \end{gathered}                                                \tag{SP18a}
\]
where $h_K$ is the class number.  These are precisely the
original conditions \cite{support-spectrum-orbit-source}.
The prime-to-$\ell$ order in (SP18a) implies that
$\psi_\ell$ has conductor dividing $\ell$; there is no
wild $\ell$-power conductor left in that factor.

\begin{corollary}[The original good-character orbit bound]
\label{cor:support-spectrum-orbit}
Fix $K$.  There are integers $C_K,A_K>0$ such that for every
$e,N\ge1$, subgroup $H\le(\mathbf Z/N\mathbf Z)^\times$,
and prime $p\nmid30Nd_K$, there is a nonzero integer $B_p$ with
\[
 \log|B_p|\le C_K\varphi([(\mathbf Z/N\mathbf Z)^\times:H])
                      [N(e+1)]^{A_K}\log p.                   \tag{SP18}
\]
For every coefficient prime $\ell\ne p$ and every original
coefficient ring $O$ for which a character $\psi$ of kernel $H$
satisfies all the original good-data hypotheses, all the
characters in that orbit have finite strict Selmer length and
\[
 \sum_{\ker\psi=H}
     \frac{\operatorname{length}_O\mathcal S(O,e,\psi)}
          {e(O/\mathbf Z_\ell)}
                    \le2v_\ell(B_p).                         \tag{SP19}
\]
Both the Selmer group and its strict local conditions are the
original ones in Corollary~\ref{cor:support-one-character-trace}.
For two distinct allowed primes $p,p'$ and any finite set $L$
of coefficient primes satisfying those hypotheses, one has
\[
 \sum_{\ell\in L}\log\ell\sum_{\ker\psi=H}
     \frac{\operatorname{length}_{O_\ell}\mathcal S(O_\ell,e,\psi)}
          {e(O_\ell/\mathbf Z_\ell)}
 \le2C_K\varphi([(\mathbf Z/N\mathbf Z)^\times:H])
          [N(e+1)]^{A_K}(\log p+\log p').                      \tag{SP20}
\]
These are the full original fixed-$K$ orbit and two-prime
requirements.  No assertion here is uniform as $K$ varies.
\end{corollary}

\begin{proof}
If the character orbit is empty, take $B_p=1$; there are no
good data to check.  Otherwise the quotient by $H$ is
cyclic of order $m=[(\mathbf Z/N\mathbf Z)^\times:H]$.
Its faithful characters form a Galois-stable set of size
$\varphi(m)$.  Choose $w\mid p$ and in Theorem~\ref{thm:support-spectrum}
use the multiset $z_\chi=\chi(Q)^{-1}$, indexed by these complex
characters, retaining repetitions of their values at $Q$.
All such values are roots of unity since $(Q,N)=1$.
Any one faithful character valued in $O$ generates all of them
by powers prime to $m$, so a coefficient embedding identifies
this index set with the original orbit over $O$.

We recall why the good-data condition survives that identification.
For $\psi' =\psi^j$, $(j,m)=1$, the order, kernel and parity
are unchanged.  The $\ell$-part of the conductor character
has order at most two in the original good-data hypotheses;
it is unchanged by such a power.  The kernel of the
remaining prime-to-$\ell$ character is unchanged as well.
These are the character-dependent conditions in the original
definition; the conditions on $K,N,\ell,e$ are fixed.
Thus every character in the orbit has the same good-data
properties.  This is the elementary argument underlying
the original orbit lemma, including its quadratic local
factor \cite{support-spectrum-orbit-source}.

At inverse arithmetic Frobenius the toral trace is exactly
$Q^{-e}\psi(Q)^{-1}+1+Q^e\psi(Q)$, not its oppositely
twisted version.  Apply
Corollary~\ref{cor:support-one-character-trace} to each orbit
character.  Its finite family consists of pairwise
nonisomorphic representations of the exact classes above.
The separation and valuation sum (SP6), summed over the
orbit, give (SP19).  The good-data hypotheses imply
$\ell\ne2$, so every coefficient prime under consideration
is allowed in (SP6).  This proves (SP18)--(SP19).

For $\ell\ne p$, use (SP19) with $B_p$; for $\ell=p$ use
the analogous integer $B_{p'}$.  Nonnegativity permits the
common bound $2v_\ell(B_p)+2v_\ell(B_{p'})$ for every
$\ell\in L$.  Unique factorization gives
$\sum_{\ell\in L}v_\ell(B)\log\ell\le\log|B|$ for every
nonzero integer $B$.  Summing and using (SP18) proves (SP20).
If $L$ is empty, the assertion is immediate; otherwise the
good-data hypotheses force $e\ge1$.  If $K$ is not imaginary
quadratic, there are no such good data, so the corresponding
source statement for arbitrary number fields is vacuous.
\end{proof}

The constants in (SP18)--(SP20) are chosen after $K$ and before
$e,N,H,p,p'$ and all coefficient primes.  The field used to
obtain automorphy upstream need not have bounded degree; the
representations counted here have already descended to the
original $K$ or $\mathbf Q$.  The corollary addresses the
original good-data range.  It supplies neither an extension
to resonant exceptional characters nor a bound uniform in a
moving imaginary quadratic field.

\begin{corollary}[The nonresonant contribution in the original family]
\label{cor:support-spectrum-case-one}
Fix a primitive Dirichlet character $\chi$ of conductor $n_0$
and order $o$, and let $F_\chi$ be the field cut out by its
kernel.  Let $b\ge7$
be the original fixed prime with $b\equiv3\pmod4$ and
$(b,n_0)=1$.  Set $K=\mathbf Q(\sqrt{-b})$.
Let $P_\chi$ be the original fixed exceptional set of odd
primes.  In particular it contains $3$ and every odd prime
dividing $on_0b h_K$.

There are an integer $C_\chi>0$ and two fixed distinct
primes $r_1,r_2$ such that the following holds for the
original arithmetic family.  Write
\[
 E=\mathbf Q(\mu_{2^a})^+F_\chi(\sqrt{-q},\sqrt{-d}),
 \qquad N=\operatorname{lcm}(2^a,n_0)qd,\qquad a\ge2,
                                                               \tag{SP21}
\]
where $q,d\equiv3\pmod4$ are the two distinct auxiliary
primes, prime to $n_0$.  Retain its original size conditions,
and assume additionally $r_i<q,d$.  For $1\le e<t$, and
any finite set $L$ of odd coefficient primes satisfying
\[
              \ell\notin P_\chi,\qquad \ell-1\nmid2e,
                                                               \tag{SP22}
\]
all the original strict Selmer lengths are finite and
\[
 \sum_{\ell\in L}\log\ell
       \sum_{\substack{\psi\in\widehat{\operatorname{Gal}(E/\mathbf Q)}
                         (O_\ell)\\\psi(-1)=(-1)^e}}
       \frac{\operatorname{length}_{O_\ell}\mathcal S(O_\ell,e,\psi)}
            {e(O_\ell/\mathbf Z_\ell)}
                      \le (Nt)^{C_\chi}.                     \tag{SP23}
\]
Here the inner set consists of the characters valued in the
original coefficient ring $O_\ell$; it need not be a full
splitting ring.  This proves the entire original Case~I
requirement for both RH and Dirichlet GRH.
\end{corollary}

\begin{proof}
We verify every good-data hypothesis for the actual family,
then sum over kernels.  The degree of $E/\mathbf Q$ divides
$2^ao$, so (SP22) makes every occurring character have
order prime to $\ell$.  At an odd prime not dividing $n_0$,
the local character of $E$ is either unramified or comes
from one of the two quadratic fields in (SP21).  Thus its
$\ell$-component has order at most two.  The modulus is even
and the parity condition is imposed in the inner sum.
The exceptional set excludes all prime divisors of $bh_K$.
The original size conditions include $b<q,d$, so
$b\nmid N$ and $(d_K,N)=1$.  Finally
$\mathbf Q(\mu_\ell)F_{\psi^{(\ell)}}$ is unramified at $b$,
because $b\ne\ell$ and the conductor of $\psi^{(\ell)}$
divides $N$.  Since $K$ is ramified at $b$, it cannot be a
subfield of that compositum.  The last good-data condition
is precisely the second inequality in (SP22).  This proves
(SP18a), including its field exclusion.  It is also a direct
proof of the original good-data verification for this family
\cite{support-spectrum-case-one-source}.

Choose two distinct primes $r_1,r_2>30n_0b$ once and for all.
When $q,d>r_1,r_2$, neither divides $30Nd_K$.  The constants
in Corollary~\ref{cor:support-spectrum-orbit} are now fixed
before $a,q,d,N,e,t$ vary.

Put $G=(\mathbf Z/N\mathbf Z)^\times$ and group all the
characters in (SP23) by their kernels $H\le G$.  For each
kernel use (SP20) on the subset of coefficient primes where
it occurs, and enlarge the character set to the full orbit
of that kernel.  All lengths are nonnegative.  A kernel of
a character has cyclic quotient, and the number of
characters with that kernel is $\varphi([G:H])$.  Distinct
kernels give disjoint sets of complex characters, so
\[
                 \sum_H\varphi([G:H])\le|\widehat G|
                                   =|G|\le N.                \tag{SP24}
\]
Consequently the left side of (SP23) is at most
\[
 D_\chi N[N(e+1)]^{A_K},\qquad
       D_\chi=2C_K(\log r_1+\log r_2).                        \tag{SP25}
\]
The original family has $N\ge16$, and $e+1\le t$.
Choose an integer $k\ge0$ with $D_\chi\le16^k$.
Then (SP25) is at most $(Nt)^{k+1+A_K}$, proving
(SP23) with $C_\chi=k+1+A_K$.

The fixed prime $b$, the class number $h_K$, and both
auxiliary primes were chosen from the fixed target
character.  Thus no varying-field constant has been
absorbed in $C_\chi$.  For ordinary RH, $n_0=o=1$ and
$F_\chi=\mathbf Q$; the same argument applies without a
separate arithmetic construction.
\end{proof}
